\section{\texttt{gui.py}: function-by-function implementation}
The first functions build immutable lessons. CryptoLab then manages Tk widgets, captured settings, worker execution, queue delivery, timeline rendering, playback and JSON export. Only the Tk thread updates widgets; cryptographic work is performed by the worker.
\subsection{\texttt{parse\_\allowbreak{}block}}
\begin{lightcode}
def parse_block(value):
\end{lightcode}
\noindent\textbf{Source:} complete \texttt{gui.py} listing in Part 11, lines 21--25.
\begin{enumerate}
\item Decode the text field as hexadecimal bytes.
\item Require exactly sixteen bytes and return the block; malformed hex or a wrong length raises ValueError for the GUI to report.
\end{enumerate}
\subsection{\texttt{snapshot}}
\begin{lightcode}
def snapshot(number, title, explanation, state, detail="", reference=None):
\end{lightcode}
\noindent\textbf{Source:} complete \texttt{gui.py} listing in Part 11, lines 28--29.
\begin{enumerate}
\item Accept a round, title, explanation, state, optional detail and optional correct reference.
\item Convert state/reference bytes to hexadecimal strings and return a dictionary. Immutable serialized states protect the timeline from subsequent in-place cipher changes and can be exported as JSON.
\end{enumerate}
\subsection{\texttt{aes\_\allowbreak{}trace}}
\begin{lightcode}
def aes_trace(block, key, decrypt=False, fault=None):
\end{lightcode}
\noindent\textbf{Source:} complete \texttt{gui.py} listing in Part 11, lines 32--57.
\begin{enumerate}
\item Start a timeline with the input block and expand keys for teaching details. In normal mode, use an after-operation observer on the actual AES encrypt/decrypt function.
\item In fault mode, use the Piret target adapter at the specified round and map its snapshots into the same schema. Return output bytes and steps; refuse faulted decryption because that lesson models encryption faults.
\end{enumerate}
\subsection{\texttt{aes\_\allowbreak{}trace.observer}}
\begin{lightcode}
def observer(number, operation, state):
\end{lightcode}
\noindent\textbf{Source:} complete \texttt{gui.py} listing in Part 11, lines 49--53.
\begin{enumerate}
\item Ignore Before events and translate each After event into a snapshot with an operation description. For AddRoundKey, include the corresponding teaching round key in the detail text.
\end{enumerate}
\subsection{\texttt{sm4\_\allowbreak{}trace}}
\begin{lightcode}
def sm4_trace(block, key, decrypt=False):
\end{lightcode}
\noindent\textbf{Source:} complete \texttt{gui.py} listing in Part 11, lines 60--86.
\begin{enumerate}
\item Start from the input block and attach an observer to the normal SM4 engine. Record rolling-state inputs and the scalar A, tau(A), linear and new-word events.
\item Add explanatory snapshots and final word reversal, using reversed subkey indices for decryption. Return the engine output and timeline; scalar substitution stages do not yet replace the displayed rolling state.
\end{enumerate}
\subsection{\texttt{sm4\_\allowbreak{}trace.observer}}
\begin{lightcode}
def observer(index, operation, value):
\end{lightcode}
\noindent\textbf{Source:} complete \texttt{gui.py} listing in Part 11, lines 65--82.
\begin{enumerate}
\item Track State input/output, collect scalar S-box intermediates and add explanatory snapshots. Use nonlocal current to retain the rolling state until its update, and append final word reversal when the Ciphertext event arrives.
\end{enumerate}
\subsection{\texttt{positions}}
\begin{lightcode}
def positions(column):
\end{lightcode}
\noindent\textbf{Source:} complete \texttt{gui.py} listing in Part 11, lines 89--91.
\begin{enumerate}
\item Accept one AES column index.
\item Delegate to piret.output\_positions through the console module and return the final ShiftRows byte mapping used by the presentation.
\end{enumerate}
\subsection{\texttt{candidates}}
\begin{lightcode}
def candidates(correct, faulty, column):
\end{lightcode}
\noindent\textbf{Source:} complete \texttt{gui.py} listing in Part 11, lines 94--96.
\begin{enumerate}
\item Accept correct/faulty AES ciphertexts and an output group.
\item Delegate to piret.key\_candidates and return its correlated quartet set. Mathematical filtering stays in the attack module.
\end{enumerate}
\subsection{\texttt{reverse\_\allowbreak{}schedule}}
\begin{lightcode}
def reverse_schedule(last_key):
\end{lightcode}
\noindent\textbf{Source:} complete \texttt{gui.py} listing in Part 11, lines 99--100.
\begin{enumerate}
\item Accept a recovered AES K10.
\item Delegate to piret.reverse\_key\_schedule and return the master key without reading GUI key fields.
\end{enumerate}
\subsection{\texttt{recover\_\allowbreak{}key}}
\begin{lightcode}
def recover_key(correct, observations):
\end{lightcode}
\noindent\textbf{Source:} complete \texttt{gui.py} listing in Part 11, lines 103--111.
\begin{enumerate}
\item Accept a correct ciphertext and round-eight faulty-output sequence, and install a local record callback.
\item Call piret.recover\_last\_key and translate observer events into GUI history tuples. Return recovered K10 or None, the history and the number of consumed faults.
\end{enumerate}
\subsection{\texttt{recover\_\allowbreak{}key.record}}
\begin{lightcode}
def record(index, column, before, fresh, survivors):
\end{lightcode}
\noindent\textbf{Source:} complete \texttt{gui.py} listing in Part 11, lines 107--108.
\begin{enumerate}
\item Translate the round-eight solver notification into a GUI history tuple containing observation, group, fresh candidates and remaining survivors. This records counts after filtering rather than performing the filtering.
\end{enumerate}
\subsection{\texttt{piret\_\allowbreak{}lesson}}
\begin{lightcode}
def piret_lesson(block, key, row, column, mask, maximum):
\end{lightcode}
\noindent\textbf{Source:} complete \texttt{gui.py} listing in Part 11, lines 114--136.
\begin{enumerate}
\item Build clean AES and round-eight faulty timelines, align matching stages and add the DFA equation explanation. Generate distinct masks and call the round-eight recovery adapter.
\item Add candidate-history snapshots, reverse a unique K10 and verify the master key. Return a lesson report; this retained round-eight lesson is distinct from the round-nine path used by the current Fault analysis selector.
\end{enumerate}
\subsection{\texttt{pdf\_\allowbreak{}aes\_\allowbreak{}lesson}}
\begin{lightcode}
def pdf_aes_lesson(block, key, row, column, mask, maximum):
\end{lightcode}
\noindent\textbf{Source:} complete \texttt{gui.py} listing in Part 11, lines 139--177.
\begin{enumerate}
\item Build a clean AES trace and process columns beginning at the selected column. For each, inject round-nine faults, show the first fault's propagation and intersect correlated key quartets.
\item Put unique quartet bytes into K10, reverse the schedule and verify the master key. If a column remains ambiguous, return an explicit incomplete report and stop rather than completing it from the teaching key.
\end{enumerate}
\subsection{\texttt{sm4\_\allowbreak{}master\_\allowbreak{}lesson}}
\begin{lightcode}
def sm4_master_lesson(block, key, start_byte, mask, maximum, giraud_mode=False, known_bit=False):
\end{lightcode}
\noindent\textbf{Source:} complete \texttt{gui.py} listing in Part 11, lines 180--208.
\begin{enumerate}
\item Record a clean SM4 lesson and define an output-only oracle plus an observer for candidate/peeling events. Choose the controlled-byte engine or the Giraud-inspired wrapper from giraud\_mode.
\item Build snapshots from clean/faulty peeled views, verify a complete master key and append a recovery or ambiguity result. Return subkeys, history, counts and display metadata; no reference subkeys are supplied to the solver.
\end{enumerate}
\subsection{\texttt{sm4\_\allowbreak{}master\_\allowbreak{}lesson.oracle}}
\begin{lightcode}
def oracle(number, byte, value):
\end{lightcode}
\noindent\textbf{Source:} complete \texttt{gui.py} listing in Part 11, lines 187--189.
\begin{enumerate}
\item Choose the controlled-byte or single-bit target adapter from the lesson mode, capture the simulation block/key and return only the requested faulty output. Do not expose trace dictionaries to the solver.
\end{enumerate}
\subsection{\texttt{sm4\_\allowbreak{}master\_\allowbreak{}lesson.show}}
\begin{lightcode}
def show(event):
\end{lightcode}
\noindent\textbf{Source:} complete \texttt{gui.py} listing in Part 11, lines 191--193.
\begin{enumerate}
\item Convert each full-recovery event into a teaching snapshot comparing faulty\_view with clean\_view. Put raw outputs, already recovered later keys, inverse-L difference and candidate survivors in the explanation details.
\end{enumerate}
\subsection{\texttt{pdf\_\allowbreak{}sm4\_\allowbreak{}lesson}}
\begin{lightcode}
def pdf_sm4_lesson(block, key, start_byte, mask, maximum):
\end{lightcode}
\noindent\textbf{Source:} complete \texttt{gui.py} listing in Part 11, lines 211--213.
\begin{enumerate}
\item Accept the block, key and selected byte/mask/observation settings.
\item Call sm4\_master\_lesson with the default controlled-byte model. Its name retains the connection to the PDF but the recovery extends the last-round example to four target rounds.
\end{enumerate}
\subsection{\texttt{aes\_\allowbreak{}giraud\_\allowbreak{}encrypt}}
\begin{lightcode}
def aes_giraud_encrypt(plaintext, key, byte=0, mask=None):
\end{lightcode}
\noindent\textbf{Source:} complete \texttt{gui.py} listing in Part 11, lines 218--220.
\begin{enumerate}
\item Accept a target block/key and optional AES byte/one-bit mask.
\item Delegate to giraud.aes\_fault\_encrypt and return its output and final-round trace. This adapter has no independent AES or candidate implementation.
\end{enumerate}
\subsection{\texttt{giraud\_\allowbreak{}lesson}}
\begin{lightcode}
def giraud_lesson(algorithm, plaintext, key, row=0, column=0, mask=1, maximum=8, known_bit=False):
\end{lightcode}
\noindent\textbf{Source:} complete \texttt{gui.py} listing in Part 11, lines 223--281.
\begin{enumerate}
\item Validate algorithm and coordinates and generate distinct one-bit masks. Route SM4 to sm4\_master\_lesson; for AES, visit every target byte beginning at the selected coordinate and intersect candidate bytes at their observed final output positions.
\item Add first-fault stage explanations and per-observation survivor sets. Once all sixteen AES K10 bytes are unique, reverse and verify the schedule; otherwise preserve an explicit ambiguous result.
\end{enumerate}
\subsection{\texttt{CryptoLab.\_\allowbreak{}\_\allowbreak{}init\_\allowbreak{}\_\allowbreak{}}}
\begin{lightcode}
def __init__(self, root):
\end{lightcode}
\noindent\textbf{Source:} complete \texttt{gui.py} listing in Part 11, lines 285--312.
\begin{enumerate}
\item Create application state, report/index/timer fields, a thread-safe queue and Tk variables for inputs, settings and status. Configure the window theme and call build.
\item Connect selection changes, keyboard navigation, close handling and periodic queue polling. No cipher is run merely by constructing the window.
\end{enumerate}
\subsection{\texttt{CryptoLab.label}}
\begin{lightcode}
def label(self, parent, text, size=11, color=TEXT):
\end{lightcode}
\noindent\textbf{Source:} complete \texttt{gui.py} listing in Part 11, lines 314--315.
\begin{enumerate}
\item Accept a parent, text, optional size and color.
\item Construct and return a consistently styled Tk Label using the parent's background. Layout is controlled by the caller.
\end{enumerate}
\subsection{\texttt{CryptoLab.button}}
\begin{lightcode}
def button(self, parent, title, action, primary=False):
\end{lightcode}
\noindent\textbf{Source:} complete \texttt{gui.py} listing in Part 11, lines 317--318.
\begin{enumerate}
\item Accept a parent, label, callback and primary-style flag.
\item Return a styled Tk Button bound to the callback. Constructing it does not invoke the action.
\end{enumerate}
\subsection{\texttt{CryptoLab.build}}
\begin{lightcode}
def build(self):
\end{lightcode}
\noindent\textbf{Source:} complete \texttt{gui.py} listing in Part 11, lines 320--434.
\begin{enumerate}
\item Construct the sidebar controls for algorithm, lesson, key, block and fault settings, then create the explorer and result tabs.
\item Allocate the three byte matrices, details box, playback/navigation controls and export action, retaining widget handles on self. Initialize disabled navigation until a lesson is available.
\end{enumerate}
\subsection{\texttt{CryptoLab.selection\_\allowbreak{}changed}}
\begin{lightcode}
def selection_changed(self, *_):
\end{lightcode}
\noindent\textbf{Source:} complete \texttt{gui.py} listing in Part 11, lines 436--452.
\begin{enumerate}
\item Ignore selector callbacks while a computation is busy. Determine whether the chosen mode is a fault lesson, Giraud lesson or SM4 and update labels and explanatory hints.
\item Enable only relevant controls: SM4 has no AES column, and exact-bit knowledge applies only to Giraud. Keep any existing report tied to its original inputs and warn that rebuilding is needed to apply new settings.
\end{enumerate}
\subsection{\texttt{CryptoLab.load\_\allowbreak{}example}}
\begin{lightcode}
def load_example(self):
\end{lightcode}
\noindent\textbf{Source:} complete \texttt{gui.py} listing in Part 11, lines 454--460.
\begin{enumerate}
\item If not busy, choose the AES or SM4 official key/input pair.
\item Load ciphertext instead of plaintext when Decrypt is selected, then update status. This changes controls only; it does not build a new timeline automatically.
\end{enumerate}
\subsection{\texttt{CryptoLab.set\_\allowbreak{}text}}
\begin{lightcode}
def set_text(self, widget, text):
\end{lightcode}
\noindent\textbf{Source:} complete \texttt{gui.py} listing in Part 11, lines 462--466.
\begin{enumerate}
\item Temporarily enable a read-only text widget, clear its content and insert the supplied text.
\item Disable the widget again. This centralizes safe updates for details and summary boxes.
\end{enumerate}
\subsection{\texttt{CryptoLab.run}}
\begin{lightcode}
def run(self):
\end{lightcode}
\noindent\textbf{Source:} complete \texttt{gui.py} listing in Part 11, lines 468--493.
\begin{enumerate}
\item Ignore a duplicate request while busy, stop playback, parse input fields and validate attack coordinates, masks and observation limits. Report invalid input without starting a worker.
\item Capture settings, disable editing and start a daemon thread running compute. The captured arguments prevent later widget changes from silently altering the running experiment.
\end{enumerate}
\subsection{\texttt{CryptoLab.compute}}
\begin{lightcode}
def compute(self, algorithm, operation, key, block, row, column, mask, maximum, known_bit=False):
\end{lightcode}
\noindent\textbf{Source:} complete \texttt{gui.py} listing in Part 11, lines 495--512.
\begin{enumerate}
\item On the worker thread, choose the Giraud, controlled-fault or normal trace builder. For ordinary encrypt/decrypt, run the inverse and check the original block is recovered.
\item Add the original settings to the report and put either the report or an error string into the queue. Do not update Tk widgets from this worker thread.
\end{enumerate}
\subsection{\texttt{CryptoLab.poll}}
\begin{lightcode}
def poll(self):
\end{lightcode}
\noindent\textbf{Source:} complete \texttt{gui.py} listing in Part 11, lines 514--539.
\begin{enumerate}
\item On the Tk thread, try to consume a queued report or error without blocking. When a result arrives, release busy state and restore the appropriate controls.
\item Display errors or install the new report, step list, matrices and summary; enable export and refresh navigation. Schedule another poll with root.after so the event loop remains responsive.
\end{enumerate}
\subsection{\texttt{CryptoLab.render}}
\begin{lightcode}
def render(self):
\end{lightcode}
\noindent\textbf{Source:} complete \texttt{gui.py} listing in Part 11, lines 541--565.
\begin{enumerate}
\item Take the selected snapshot and choose its explicit correct reference or the previous displayed state. Decode both blocks, compute XOR and changed-byte/bit counts, and choose AES column-major or SM4 word-row layout.
\item Update titles, three matrices, red highlights, detail text and navigation indicators. Changing the displayed step never reruns encryption or changes solver results.
\end{enumerate}
\subsection{\texttt{CryptoLab.navigation}}
\begin{lightcode}
def navigation(self):
\end{lightcode}
\noindent\textbf{Source:} complete \texttt{gui.py} listing in Part 11, lines 567--573.
\begin{enumerate}
\item Determine whether a report is available and computation is idle.
\item Enable or disable Previous, Next, Play, step selection and the slider based on this state and current index. This guards navigation at the timeline boundaries.
\end{enumerate}
\subsection{\texttt{CryptoLab.select}}
\begin{lightcode}
def select(self, index):
\end{lightcode}
\noindent\textbf{Source:} complete \texttt{gui.py} listing in Part 11, lines 575--579.
\begin{enumerate}
\item If a report is ready and the application is idle, stop automatic playback.
\item Store the supplied valid step index and call render. Internal callers obtain indices from the slider, dropdown or bounded movement methods.
\end{enumerate}
\subsection{\texttt{CryptoLab.seek\_\allowbreak{}to}}
\begin{lightcode}
def seek_to(self, value):
\end{lightcode}
\noindent\textbf{Source:} complete \texttt{gui.py} listing in Part 11, lines 581--584.
\begin{enumerate}
\item Convert the slider's textual numeric value to an integer step index.
\item If it differs from the displayed index and a report exists, delegate to select. Avoiding redundant selection also prevents unnecessary re-rendering.
\end{enumerate}
\subsection{\texttt{CryptoLab.move}}
\begin{lightcode}
def move(self, direction):
\end{lightcode}
\noindent\textbf{Source:} complete \texttt{gui.py} listing in Part 11, lines 586--588.
\begin{enumerate}
\item Require an existing report and add the requested direction to the current index.
\item Clamp the result to the first and last valid indices and call select. Previous/Next and keyboard actions share this boundary handling.
\end{enumerate}
\subsection{\texttt{CryptoLab.play}}
\begin{lightcode}
def play(self):
\end{lightcode}
\noindent\textbf{Source:} complete \texttt{gui.py} listing in Part 11, lines 590--600.
\begin{enumerate}
\item If a timer is already scheduled, stop playback; otherwise require an idle application and report. Restart from the first step when already at the end.
\item Set the Pause label, render and schedule tick after the configured millisecond delay. Tk's event loop remains in control rather than being blocked by a sleep loop.
\end{enumerate}
\subsection{\texttt{CryptoLab.tick}}
\begin{lightcode}
def tick(self):
\end{lightcode}
\noindent\textbf{Source:} complete \texttt{gui.py} listing in Part 11, lines 602--609.
\begin{enumerate}
\item Clear the stored timer identifier and advance by one timeline step.
\item Render the new step, schedule another tick when more steps remain, or stop at the end. This method advances a precomputed lesson rather than executing a cipher round.
\end{enumerate}
\subsection{\texttt{CryptoLab.stop}}
\begin{lightcode}
def stop(self):
\end{lightcode}
\noindent\textbf{Source:} complete \texttt{gui.py} listing in Part 11, lines 611--615.
\begin{enumerate}
\item Cancel a scheduled after callback if one exists and clear its identifier.
\item Restore the Play label. It is safe to call this while no timer is active.
\end{enumerate}
\subsection{\texttt{CryptoLab.fill\_\allowbreak{}summary}}
\begin{lightcode}
def fill_summary(self):
\end{lightcode}
\noindent\textbf{Source:} complete \texttt{gui.py} listing in Part 11, lines 617--633.
\begin{enumerate}
\item Read the installed report and assemble input/output, model and verification information. For fault lessons, include used observations, recovered key/subkeys and ambiguity or bit-knowledge details.
\item For normal execution, include inverse verification and recovered input. Write the text through set\_text, explicitly distinguishing simulation data and teaching-key export from a physical experiment.
\end{enumerate}
\subsection{\texttt{CryptoLab.export}}
\begin{lightcode}
def export(self):
\end{lightcode}
\noindent\textbf{Source:} complete \texttt{gui.py} listing in Part 11, lines 635--645.
\begin{enumerate}
\item Require a report and ask the user for a JSON filename.
\item Serialize the complete current report with indentation and update status; catch file-writing errors in a dialog. The export includes the teaching key and hidden snapshots, so it is a lesson archive rather than an attacker-only observation file.
\end{enumerate}
\subsection{\texttt{CryptoLab.close}}
\begin{lightcode}
def close(self):
\end{lightcode}
\noindent\textbf{Source:} complete \texttt{gui.py} listing in Part 11, lines 647--650.
\begin{enumerate}
\item Stop playback and cancel periodic queue polling.
\item Destroy the root window. No cryptographic key recovery is performed during shutdown.
\end{enumerate}

\subsection{Anonymous event handlers and the worker lifecycle}
The lambdas created in \texttt{build} connect buttons to movement methods, preserving their direction argument. Keyboard lambdas suppress navigation when focus is in a text entry or dropdown. The step-selector handler reads the chosen list index; the speed control supplies a millisecond delay, not the number of cipher rounds. These expressions return presentation actions rather than candidate bytes.

The complete execution order after clicking Build lesson is:
\[
\texttt{run} \longrightarrow \texttt{compute (worker)} \longrightarrow
\texttt{messages queue} \longrightarrow \texttt{poll (Tk thread)}.
\]
After installing the report, \texttt{render} displays a snapshot and \texttt{fill\_\allowbreak{}summary} displays the result. Playback follows \texttt{play}, repeated \texttt{tick}, and \texttt{stop}; export follows \texttt{export} and writes this already computed report. Neither playback nor export repeats a fault experiment.

\section{Return values, mutation and error handling at a glance}
\begin{description}
\item[Normal cipher outputs:] AES and SM4 block functions return sixteen immutable bytes. They do not print, pause, add message padding, choose an encryption mode or authenticate a message.
\item[Trace outputs:] target adapters return ciphertext and a trace; GUI trace functions return output and snapshot dictionaries. These teacher views may contain internal states. Solver calls receive a deliberately smaller interface.
\item[Candidate outputs:] AES Piret returns a set of four-byte tuples; AES Giraud and SM4 filters return byte candidate sets and position/model diagnostics. A tuple correlation cannot be replaced by choosing four independent byte survivors.
\item[Recovery outputs:] complete keys are returned only after uniqueness. SM4 reports include \texttt{complete}, \texttt{master\_\allowbreak{}key}, \texttt{round\_\allowbreak{}keys}, \texttt{history}, \texttt{candidate\_\allowbreak{}counts} and \texttt{used}. Incomplete recovery returns no master key, even if some subkeys are already known.
\item[Errors:] invalid sizes/masks and incompatible differential support raise \texttt{ValueError}; failed verification raises \texttt{RuntimeError}. A compatible but nonunique intersection is reported as ambiguity rather than being mislabeled an exception or a successful key guess.
\end{description}

\section{Worked reading exercises}
\begin{enumerate}
\item Follow one AES byte from input index one through \texttt{bytes\_\allowbreak{}to\_\allowbreak{}state}, initial AddRoundKey, SubBytes and ShiftRows. Explain why its flat output index changes even though SubBytes itself does not move bytes.
\item Compare \texttt{gf\_\allowbreak{}mul} with integer multiplication and \texttt{linear} with \texttt{linear\_\allowbreak{}key}. Identify the exact call sites that would become wrong if either pair were interchanged.
\item At one round-nine Piret observation, read the four tables in \texttt{key\_\allowbreak{}candidates}. Check that every surviving quartet shares one error row and one error magnitude; explain why four separately plausible bytes do not necessarily form a valid quartet.
\item Run a Giraud byte pair through \texttt{aes\_\allowbreak{}candidates} twice: with \texttt{known\_\allowbreak{}mask=None} and with the exact mask. Identify the allowed difference set in each case and why the stronger filter can leave fewer candidates.
\item In SM4, follow the first \texttt{recover\_\allowbreak{}sm4\_\allowbreak{}master\_\allowbreak{}key} event for round thirty and inspect \texttt{peeled\_\allowbreak{}with}. It must contain recovered round thirty-one, not reference round thirty or the original master key.
\item In the GUI, compare an ordinary transformation snapshot with a fault snapshot. Explain why changing from the previous state is not the same measurement as a faulty-versus-correct state difference, even though both use an XOR and red highlights.
\end{enumerate}
